\section{Three exact non-descent theorems}
\label{sec:non-descent}

The three failures below use different data. The clock theorem does not
depend on the marker convention. The marker/multiplicity theorem does not use
the analytic coefficient limit. The determinant theorem does not require an
objectwise map. Keeping them separate shows exactly which proposed
identification each witness defeats.

\subsection{Clock preservation forces composite support}

\begin{theorem}[No total exact-clock rational-prime map]
\label{thm:clock}
There is no total map
\[
 \pi:\Prim_q\longrightarrow\Primes
\]
such that
\[
 \log\pi(\gamma)=n(\gamma)\log q
\]
for every primitive source necklace $\gamma$.
\end{theorem}

\begin{proof}
The word $01$ exists over each frozen alphabet. It is primitive: if a
two-letter word were a proper power, it would be the square of a one-letter
word and its two symbols would coincide. Let $\gamma=[01]$. Clock
preservation gives
\[
 \log\pi(\gamma)=2\log q=\log(q^2).
\]
The real logarithm is injective on positive numbers, so
$\pi(\gamma)=q^2$. Since $q>1$, this forced label is composite, contrary to
$\pi(\gamma)\in\Primes$.
\end{proof}

Totality matters because the claim concerns the full primitive product.
Omitting $[01]$ defines a partial projection, not a counterexample to
\cref{thm:clock}. Mapping the class to a finite-field prime polynomial also
changes the target type, and replacing the clock by an arbitrary prime
enumeration abandons the theorem's equality.

\subsection{Marker and weight expose multiplicity}

\begin{theorem}[No marker--weight--multiplicity identification]
\label{thm:multiplicity}
There is no factorwise identification of all source primitive factors
\[
 \left(1-z^nq^{-ns}\right)^{-1}
\]
with rational-prime factors $(1-zp^{-s})^{-1}$ that preserves the free
marker, analytic weight, and multiplicity.
\end{theorem}

\begin{proof}
Equality of the factor monomials requires
\begin{equation}
 z^nq^{-ns}=zp^{-s}.                                         \label{eq:factor-equality}
\end{equation}
Because $z$ is free, equality of exponents in \cref{eq:factor-equality}
forces $n=1$. Equality of the weights as functions of $s$ then forces $p=q$.
By \cref{eq:necklace},
\[
 \Nq(1)=q.
\]
Thus the source has $q$ distinct length-one primitive necklaces, all with the
only compatible monomial $zq^{-s}$. The rational Euler product has one
primitive factor indexed by the rational prime $p=q$. The required factor
multiplicity is therefore $q:1$, not $1:1$.
\end{proof}

Specializing $z=1$ before the comparison erases the marker witness rather
than answering it. Likewise, merging the $q$ source factors or selecting one
of them introduces an explicit quotient or projection. Those constructions
can be useful controls, but they are outside the factorwise identification in
\cref{thm:multiplicity}.

\begin{figure}[t]
  \centering
  \resizebox{\textwidth}{!}{\begin{tikzpicture}[
  font=\small,
  >=Latex,
  channel/.style={draw=charcoal, rounded corners=2pt, very thick,
    align=center, text width=43mm, minimum height=13mm, fill=softgray},
  premise/.style={channel, draw=formalblue, fill=formalblue!8},
  witness/.style={channel, draw=governancepurple, fill=governancepurple!8},
  stop/.style={channel, draw=warningamber, fill=warningamber!9},
  flow/.style={->, very thick, draw=charcoal},
  fail/.style={->, very thick, dashed, draw=warningamber}
]
\node[premise] (c1) at (0,0)
  {clock/support\\total $\pi$ and $\log p=n\log q$};
\node[witness] (w1) at (0,-1.8)
  {$[01]$ primitive, $n=2$\\forced label $p=q^2$};
\node[stop] (s1) at (0,-3.6)
  {$q^2$ composite\\no rational-prime support};

\node[premise] (c2) at (5.1,0)
  {marker/multiplicity\\$z^nq^{-ns}=zp^{-s}$};
\node[witness] (w2) at (5.1,-1.8)
  {$n=1$, $p=q$\\$N_q(1)=q$};
\node[stop] (s2) at (5.1,-3.6)
  {$q$ source factors, one target\\multiplicity $q:1$};

\node[premise] (c3) at (10.2,0)
  {complete determinants\\compare $[z](-\log D)$};
\node[witness] (w3) at (10.2,-1.8)
  {source $q^{1-s}$\\target $P(s)=\sum_pp^{-s}$};
\node[stop] (s3) at (10.2,-3.6)
  {scaled limits differ\\analytic nonidentity};

\draw[flow] (c1) -- (w1);
\draw[fail] (w1) -- (s1);
\draw[flow] (c2) -- (w2);
\draw[fail] (w2) -- (s2);
\draw[flow] (c3) -- (w3);
\draw[fail] (w3) -- (s3);

\node[draw=charcoal, rounded corners=2pt, align=center, fill=white,
  text width=138mm, below=7mm of s2]
  {Independent channels: the clock proof uses no marker; the multiplicity
   proof uses no coefficient limit; the determinant proof uses no factor map.};
\end{tikzpicture}
}
  \caption{Three independent exact failure channels. Clock preservation sends
  the primitive class $[01]$ to the composite $q^2$. Marker and weight
  preservation reduce to length one, where $q$ source factors collide with
  one target factor. The complete marked determinants already differ in the
  first logarithmic $z$ coefficient.}
  \label{fig:three-obstructions}
\end{figure}

\subsection{The first marked coefficient already differs}

\begin{theorem}[Marked determinant nonidentity]
\label{thm:coefficient}
On their common absolute-convergence domain, $\Dq(s,z)$ and $\DP(s,z)$ are
not equal as analytic functions of $(s,z)$. Their logarithms differ in the
coefficient of $z$.
\end{theorem}

\begin{proof}
Equations \eqref{eq:source-det} and \eqref{eq:target-log} give
\begin{equation}
 [z]\{-\log\Dq(s,z)\}=q^{1-s},
 \qquad
 [z]\{-\log\DP(s,z)\}=P(s)=\sum_{p\in\Primes}p^{-s}.         \label{eq:first-coefficients}
\end{equation}
The target series is the coefficient obtained from the rational Euler product
on $\RePart s>1$ \citep{nistDLMF25211}. We show the two functions in
\cref{eq:first-coefficients} cannot agree.

Set $s=\sigma>1$ real and multiply by $2^\sigma$. For $q=2$, the source is
identically $2$, whereas
\[
 2^\sigma P(\sigma)
 =1+\sum_{p\ge3}\left(\frac2p\right)^\sigma\longrightarrow1.
\]
For $q=3$ or $q=5$, the scaled source is
\[
 2^\sigma q^{1-\sigma}=q\left(\frac2q\right)^\sigma
 \longrightarrow0,
\]
while the scaled target still tends to one. To justify the target limit,
bound the tail for $\sigma\ge2$ by
$\sum_{n\ge3}(2/n)^2<\infty$ and apply dominated convergence term by term.
Thus the first coefficients differ for every frozen $q$, and the analytic
determinants cannot be identical.
\end{proof}

This argument uses no rational-prime table, Riemann-zero data, fitting, or
global divisor comparison. It also has the correct determinant orientation:
the positive Euler products are reciprocals of $\Dq$ and $\DP$, while
\cref{thm:coefficient} compares $-\log$ of the determinants.

\begin{corollary}[Three bounded closures]
\label{cor:closures}
The total exact-clock rational-prime projection fails; the same-marker
factorwise identification fails; and equality of the complete marked
determinants fails. No one conclusion is enlarged to all partial maps,
changed markers, induced systems, or symbolic extensions.
\end{corollary}

\begin{proof}
Apply \cref{thm:clock,thm:multiplicity,thm:coefficient} to their respective
contracts. Their quantifier notes exclude precisely the changed constructions
listed in the statement.
\end{proof}
